\chapter{Literature Review} \label{chap:ch2}

\section{Introduction}\label{sec:se21}

This literature review followed \textcite{Kitchenham2009}'s 
Systematic Literature Review workflow.
% as outlined by Kitchenham and Charters (2007)~\cite{kitchenham2007guidelines}.
The review aimed to identify and analyze existing research relevant to the
subject matter of this thesis. Likewise, it was designed to provide
a strong foundation for the following topics:

\begin{enumerate}
    \item Model Partition and Split Inference
    \item \ac{ML} Orchestration/Scheduling in Edge Computing
    \item Computational Offloading for \ac{ML} inference
    \item Heterogeneous Resource Management in \ac{IoT}/Edge
\end{enumerate}

In addition, the proposal's bibliography was also
included and examined in this context of this work.


\section{Review Methodology}\label{sec:se22}
\subsection{Search Strategy} \label{sec:se221}

There were three digital libraries used to search for relevant articles:
the IEEE Xplore Digital Library \cite{IEEEXplore} and the ACM Digital
Library \cite{ACMDigitalLibrary} as the main sources of specialized
literature on the subject, and Scopus \cite{Scopus} as a broader database
to cover any gaps that could arise.

To accomplish this search, the following search strings were built,
respectively for each topic:\footnote{It should be noted that the
proximity and wildcard operators were adapted to each digital library's syntax.}
\begin{enumerate}
    \item ("split inference" OR "distributed inference" OR "split computing" OR "collaborative inference" OR "layer partitioning" OR "DNN partitioning") AND ("Edge" OR "Far-Edge" OR "TinyML" OR "microcontroller" OR "resource-constrained")	
    \item ( ("resource-aware" OR "network-aware" OR "adaptive") NEAR/5 ("schedul*" OR "orchestrat*" OR "placement" OR "allocation") ) AND ( "inference" OR "ML model*" OR "DNN" ) AND ( "Edge" OR "Fog" OR "IoT" OR "Far-Edge" )
    \item ( "inference" NEAR/3 ("offload*" OR "decision*") OR "inference offloading" ) AND ( "resource-aware" OR "latency-aware" OR "energy-aware" OR "adaptive" OR "resource-constrained" OR "decision model" ) AND ( "Edge" OR "Far-Edge" OR "IoT" OR "TinyML" OR "MCU" )	
    \item ("device heterogeneity" OR "hardware heterogeneity" OR "heterogeneous computing" OR "heterogeneous resources" OR "heterogeneous cluster") AND ("schedul*" OR "orchestrat*" OR "allocation") AND ("inference" OR "DNN" OR "ML workload*") AND ("Edge" OR "IoT" OR "Far-Edge")	
\end{enumerate}

\subsection{Inclusion and Exclusion Criteria} \label{sec:se222}

In order to filter the search results, the rules in Table \ref{tab:lr_incl} were applied to include or
exclude the documents found:

\begin{table}[ht]
\centering
\begin{tabular}{|c|>{\raggedright\arraybackslash}p{0.5\textwidth}|}
\hline
1 & Results found during the search \\
\hline
2 & Date of publication since 2018 (inclusive) \\
\hline
3 & Full-text available \\
\hline
4 & Published in the English language \\
\hline
5 & Adequate document format\\
\hline
6 & Not duplicated \\
\hline
... & ... \\
\hline

\end{tabular}
\vspace{2mm}
\caption{Inclusion and exclusion criteria for the systematic literature review.}\label{tab:lr_incl}
\end{table}

Different formats were found in the indicated set of digital libraries.
The following formats were considered adequate for this review: 
\emph{Conference Paper}, \emph{Journal} and \emph{Magazine} for \emph{IEEE Xplore};
\emph{Research-Article}, \emph{Demonstration} and \emph{Short-Paper},
%\emph{Poster} and \emph{Survey} 
for \emph{ACM Digital Library};
\emph{Article}, \emph{Conference Paper} and \emph{Conference Review},
\emph{Book Chapter}, \emph{Review} and \emph{Book} for \emph{Scopus}.
The document type filter also excluded \emph{Work-in-Progress}
documents.

\section{Results of the Search} \label{sec:se23}

The literature search was fruitful, returning a total of 1,577
documents across the four topics and the three digital libraries used.

Table \ref{tab:lr_lib} summarizes the number of documents found on the initial
stage per topic and digital library, yet to be filtered through the inclusion
and exclusion criteria.

\begin{table}[ht]
\centering
\begin{tabular}{|c|c|c|c|c|}
\hline
\textbf{Topic} & \textbf{IEEE Xplore} & \textbf{ACM DL} & \textbf{Scopus} & \textbf{Total} \\
\hline
1 & 455 &  19 & 587 & 1061 \\
\hline
2 & 91 & 1 & 97 & 189 \\
\hline
3 & 68 & 22 & 96 & 186 \\
\hline
4 & 88 & 0 & 53 & 141 \\
\hline
\textbf{Total} & 702 & 42 & 833 & \textbf{1577} \\
\hline
\end{tabular}
\vspace{2mm}
\caption{Number of documents found per topic and digital library.}\label{tab:lr_lib}
\end{table}

After applying the inclusion and exclusion criteria, a total of ?? documents
were included for further analysis and serve as basis for this work.
Figure \ref{fig:lr_screening} illustrates the flow that documents went through
during the screening process. 

\begin{figure}
\centering
    \includegraphics[width=0.8\textwidth]{Literature_review_screening.png}
\caption{Screening process of the systematic literature review.}\label{fig:lr_screening}
\end{figure}